% ============================================================
% PREGUNTAS POST-CLASE: ELECTROMAGNETISMO
% Curso de Oriente - CCH Oriente | Enero 2026
% Solo contenido — pegar en plantilla Overleaf
% ============================================================

\renewcommand{\tema}{Preguntas Post-Clase — Electromagnetismo}

\twocolumn
\section{PREGUNTAS POST-CLASE}

% ============================================================
% PREGUNTAS 1–5: Aplicación directa de conceptos
% ============================================================

\begin{preguntabox}
\subsection{Pregunta 1}
Dos electrodos de una celda de combustible tienen cargas de $4 \times 10^{-6}$ C y $9 \times 10^{-6}$ C, separados 0.6 m. ¿Cuál es la magnitud de la fuerza eléctrica entre ellos? ($k = 9 \times 10^9$ N$\cdot$m$^2$/C$^2$)

\begin{enumerate}
    \item[A)] 0.6 N
    \item[B)] 0.9 N
    \item[C)] 1.8 N
    \item[D)] 9 N
\end{enumerate}
\end{preguntabox}

\begin{preguntabox}
\subsection{Pregunta 2}
La antena de una estación base de telefonía celular porta una carga de $5 \times 10^{-6}$ C. ¿Cuál es la magnitud del campo eléctrico a 3 m de la antena? ($k = 9 \times 10^9$ N$\cdot$m$^2$/C$^2$)

\begin{enumerate}
    \item[A)] 1500 N/C
    \item[B)] 3000 N/C
    \item[C)] 5000 N/C
    \item[D)] 15000 N/C
\end{enumerate}
\end{preguntabox}

\begin{preguntabox}
\subsection{Pregunta 3}
El driver de un foco LED tiene una resistencia de 3 $\Omega$ y opera con una batería de 9 V. ¿Qué corriente circula por el circuito?

\begin{enumerate}
    \item[A)] 0.33 A
    \item[B)] 1 A
    \item[C)] 3 A
    \item[D)] 27 A
\end{enumerate}
\end{preguntabox}

\begin{preguntabox}
\subsection{Pregunta 4}
El motor de un dron consume una corriente de 5 A y tiene una resistencia interna de 8 $\Omega$. ¿Cuál es la potencia disipada en la resistencia?

\begin{enumerate}
    \item[A)] 40 W
    \item[B)] 65 W
    \item[C)] 160 W
    \item[D)] 200 W
\end{enumerate}
\end{preguntabox}

\begin{preguntabox}
\subsection{Pregunta 5}
Tres resistencias de 10 $\Omega$, 15 $\Omega$ y 5 $\Omega$ se conectan en serie al sistema de iluminación de un semáforo (90 V). ¿Qué corriente circula por el circuito?

\begin{enumerate}
    \item[A)] 1 A
    \item[B)] 2 A
    \item[C)] 3 A
    \item[D)] 9 A
\end{enumerate}
\end{preguntabox}

% ============================================================
% PREGUNTAS 6–10: Combinación de conceptos
% ============================================================

\begin{preguntabox}
\subsection{Pregunta 6}
Dos motores eléctricos de una cinta transportadora tienen resistencias de 20 $\Omega$ y 30 $\Omega$ conectadas en paralelo a 60 V. ¿Cuál es la corriente total que suministra la fuente?

\begin{enumerate}
    \item[A)] 2 A
    \item[B)] 3 A
    \item[C)] 5 A
    \item[D)] 12 A
\end{enumerate}
\end{preguntabox}

\begin{preguntabox}
\subsection{Pregunta 7}
Un sistema de audio de alta fidelidad utiliza tres condensadores de 3 $\mu$F, 5 $\mu$F y 2 $\mu$F conectados en paralelo como filtro de frecuencias. ¿Cuál es la capacitancia equivalente del conjunto?

\begin{enumerate}
    \item[A)] 1 $\mu$F
    \item[B)] 3 $\mu$F
    \item[C)] 8 $\mu$F
    \item[D)] 10 $\mu$F
\end{enumerate}
\end{preguntabox}

\begin{preguntabox}
\subsection{Pregunta 8}
El condensador de un equipo de desfibrilación cardíaca tiene una capacitancia de 100 $\mu$F y se carga a 200 V. ¿Cuánta energía almacena?

\begin{enumerate}
    \item[A)] 0.5 J
    \item[B)] 1 J
    \item[C)] 2 J
    \item[D)] 20 J
\end{enumerate}
\end{preguntabox}

\begin{preguntabox}
\subsection{Pregunta 9}
Un conductor recto de alta tensión conduce una corriente de 10 A. ¿Cuál es la magnitud del campo magnético a 0.1 m del conductor? ($\mu_0 = 4\pi \times 10^{-7}$ T$\cdot$m/A)

\begin{enumerate}
    \item[A)] $10 \times 10^{-6}$ T
    \item[B)] $20 \times 10^{-6}$ T
    \item[C)] $40 \times 10^{-6}$ T
    \item[D)] $200 \times 10^{-6}$ T
\end{enumerate}
\end{preguntabox}

\begin{preguntabox}
\subsection{Pregunta 10}
Un aerogenerador produce una fem inducida gracias al movimiento de sus bobinas en un campo magnético. Si una bobina de 200 vueltas experimenta un cambio de flujo de 0.05 Wb en 0.1 s, ¿cuál es la fem inducida?

\begin{enumerate}
    \item[A)] 10 V
    \item[B)] 50 V
    \item[C)] 100 V
    \item[D)] 1000 V
\end{enumerate}
\end{preguntabox}

% ============================================================
% PREGUNTAS 11–15: Problemas desafiantes con múltiples pasos
% ============================================================

\begin{preguntabox}
\subsection{Pregunta 11}
En la instalación eléctrica de un laboratorio, una resistencia de 6 $\Omega$ está conectada en serie con dos resistencias de 12 $\Omega$ cada una, las cuales están en paralelo entre sí. El conjunto se conecta a 72 V. ¿Cuál es la corriente total que sale de la fuente?

\begin{enumerate}
    \item[A)] 3 A
    \item[B)] 4 A
    \item[C)] 6 A
    \item[D)] 12 A
\end{enumerate}
\end{preguntabox}

\begin{preguntabox}
\subsection{Pregunta 12}
En el circuito del problema anterior (6 $\Omega$ en serie con dos ramas de 12 $\Omega$ en paralelo, conectado a 72 V), ¿qué corriente circula por \textit{cada una} de las ramas de 12 $\Omega$?

\begin{enumerate}
    \item[A)] 1.5 A
    \item[B)] 2 A
    \item[C)] 3 A
    \item[D)] 6 A
\end{enumerate}
\end{preguntabox}

\begin{preguntabox}
\subsection{Pregunta 13}
En el sistema de arranque de un vehículo híbrido, un condensador de 6 $\mu$F está conectado en serie con un grupo de dos condensadores de 4 $\mu$F y 8 $\mu$F en paralelo. ¿Cuál es la capacitancia equivalente total del sistema?

\begin{enumerate}
    \item[A)] 2 $\mu$F
    \item[B)] 4 $\mu$F
    \item[C)] 6 $\mu$F
    \item[D)] 18 $\mu$F
\end{enumerate}
\end{preguntabox}

\begin{preguntabox}
\subsection{Pregunta 14}
Una espira cuadrada de 0.4 m$^2$ de área se encuentra en un campo magnético uniforme de 0.5 T. El plano de la espira forma un ángulo de 30$^\circ$ con el campo magnético. ¿Cuál es el flujo magnético a través de la espira?

\begin{enumerate}
    \item[A)] 0.1 Wb
    \item[B)] 0.17 Wb
    \item[C)] 0.2 Wb
    \item[D)] 0.35 Wb
\end{enumerate}
\end{preguntabox}

\begin{preguntabox}
\subsection{Pregunta 15}
El transformador de una subestación eléctrica induce una fem de 300 V en una bobina. Si el cambio de flujo magnético es de 0.06 Wb y ocurre en 0.04 s, ¿cuántas vueltas tiene la bobina?

\begin{enumerate}
    \item[A)] 50 vueltas
    \item[B)] 100 vueltas
    \item[C)] 200 vueltas
    \item[D)] 450 vueltas
\end{enumerate}
\end{preguntabox}

% ============================================================
\newpage
\onecolumn
\section{RESPUESTAS Y RETROALIMENTACIÓN}
% ============================================================

\begin{respuestabox}
\subsection{Pregunta 1}
\textcolor{verdecorrecto}{\textbf{Respuesta correcta: B) 0.9 N}}

\textbf{Justificación:}\\
\textbf{Fórmula utilizada:} Ley de Coulomb (Lección: Sección 1 — Electrostática y Ley de Coulomb):
$$F = k \frac{|q_1| \cdot |q_2|}{r^2}$$

\textbf{Sustitución:}
$$F = (9 \times 10^9 \text{ N}\cdot\text{m}^2/\text{C}^2) \cdot \frac{(4 \times 10^{-6} \text{ C})(9 \times 10^{-6} \text{ C})}{(0.6 \text{ m})^2}$$
$$F = (9 \times 10^9) \cdot \frac{36 \times 10^{-12} \text{ C}^2}{0.36 \text{ m}^2} = (9 \times 10^9) \cdot (10^{-10}) = 0.9 \text{ N}$$

\textbf{Respuestas incorrectas:}
\begin{itemize}
    \item \textbf{A) 0.6 N:} Error por usar $r$ en lugar de $r^2$ en el denominador: $k \cdot q_1 q_2 / r$ daría otro valor; posiblemente confusión aritmética.
    \item \textbf{C) 1.8 N:} Error por no elevar $r$ al cuadrado: se usó $r^2 = 0.36$ pero se duplicó el resultado sin justificación.
    \item \textbf{D) 9 N:} Error por omitir el denominador $r^2$, calculando solo $k \cdot q_1 \cdot q_2 \times 10$ — no se aplicó la ley de inverso del cuadrado.
\end{itemize}
\end{respuestabox}

\begin{respuestabox}
\subsection{Pregunta 2}
\textcolor{verdecorrecto}{\textbf{Respuesta correcta: C) 5000 N/C}}

\textbf{Justificación:}\\
\textbf{Fórmula utilizada:} Campo eléctrico de una carga puntual (Lección: Sección 1 — Campo eléctrico):
$$E = k \frac{|q|}{r^2}$$

\textbf{Sustitución:}
$$E = (9 \times 10^9 \text{ N}\cdot\text{m}^2/\text{C}^2) \cdot \frac{5 \times 10^{-6} \text{ C}}{(3 \text{ m})^2} = \frac{9 \times 10^9 \times 5 \times 10^{-6}}{9 \text{ m}^2} = \frac{45000}{9} = 5000 \text{ N/C}$$

\textbf{Respuestas incorrectas:}
\begin{itemize}
    \item \textbf{A) 1500 N/C:} Error por dividir entre $r^2 = 9$ y luego dividir de nuevo entre 3, como si hubiera dos factores de distancia.
    \item \textbf{B) 3000 N/C:} Error por usar $q = 3 \times 10^{-6}$ C en lugar de $5 \times 10^{-6}$ C — confusión con la distancia $r = 3$ m.
    \item \textbf{D) 15000 N/C:} Error por usar $r$ en lugar de $r^2$ en el denominador: $k \cdot q / r = 9\times10^9 \times 5\times10^{-6} / 3 = 15000$.
\end{itemize}
\end{respuestabox}

\begin{respuestabox}
\subsection{Pregunta 3}
\textcolor{verdecorrecto}{\textbf{Respuesta correcta: C) 3 A}}

\textbf{Justificación:}\\
\textbf{Fórmula utilizada:} Ley de Ohm despejada para corriente (Lección: Sección 2 — Resistencia eléctrica y Ley de Ohm):
$$I = \frac{V}{R}$$

\textbf{Sustitución:}
$$I = \frac{9 \text{ V}}{3 \text{ }\Omega} = 3 \text{ A}$$

\textbf{Respuestas incorrectas:}
\begin{itemize}
    \item \textbf{A) 0.33 A:} Error por invertir el cociente: $R/V = 3/9 = 0.33$ — se dividió la resistencia entre el voltaje.
    \item \textbf{B) 1 A:} Error por confundir la corriente con la diferencia $V - R = 9 - 3$ dividida entre algo, o por usar una fórmula incorrecta.
    \item \textbf{D) 27 A:} Error por multiplicar $V \times R = 9 \times 3 = 27$ en lugar de dividir — inversión conceptual de la Ley de Ohm.
\end{itemize}
\end{respuestabox}

\begin{respuestabox}
\subsection{Pregunta 4}
\textcolor{verdecorrecto}{\textbf{Respuesta correcta: D) 200 W}}

\textbf{Justificación:}\\
\textbf{Fórmula utilizada:} Potencia en función de corriente y resistencia (Lección: Sección 2 — Potencia eléctrica):
$$P = I^2 \cdot R$$

\textbf{Sustitución:}
$$P = (5 \text{ A})^2 \times 8 \text{ }\Omega = 25 \text{ A}^2 \times 8 \text{ }\Omega = 200 \text{ W}$$

\textbf{Respuestas incorrectas:}
\begin{itemize}
    \item \textbf{A) 40 W:} Error por usar $P = I \times R$ sin elevar la corriente al cuadrado: $5 \times 8 = 40$.
    \item \textbf{B) 65 W:} Error aritmético al calcular $I^2 = 25$ pero usar $25 + 8 \times 5 = 25 + 40 = 65$ — mezcla de operaciones incorrectas.
    \item \textbf{C) 160 W:} Error por calcular $I^2 = 25$ y luego usar $R = 6.4\,\Omega$ en lugar de $8\,\Omega$, o por hacer $5 \times 8 \times 4$ con un factor adicional erróneo.
\end{itemize}
\end{respuestabox}

\begin{respuestabox}
\subsection{Pregunta 5}
\textcolor{verdecorrecto}{\textbf{Respuesta correcta: C) 3 A}}

\textbf{Justificación:}\\
\textbf{Fórmula utilizada:} Resistencia equivalente en serie y Ley de Ohm (Lección: Sección 3 — Resistencias en serie):
$$R_{eq} = R_1 + R_2 + R_3 \qquad I = \frac{V}{R_{eq}}$$

\textbf{Sustitución:}
$$R_{eq} = 10 \text{ }\Omega + 15 \text{ }\Omega + 5 \text{ }\Omega = 30 \text{ }\Omega$$
$$I = \frac{90 \text{ V}}{30 \text{ }\Omega} = 3 \text{ A}$$

\textbf{Respuestas incorrectas:}
\begin{itemize}
    \item \textbf{A) 1 A:} Error por dividir el voltaje entre la resistencia mayor únicamente: $90/15 \neq 1$.
    \item \textbf{B) 2 A:} Error por sumar solo dos resistencias: $R_{eq} = 10 + 15 = 25$... o por usar $R_{eq} = 45\,\Omega$ con un error de suma.
    \item \textbf{D) 9 A:} Error por dividir el voltaje entre la resistencia menor: $90/10 = 9$ — se ignoran las otras dos resistencias.
\end{itemize}
\end{respuestabox}

\begin{respuestabox}
\subsection{Pregunta 6}
\textcolor{verdecorrecto}{\textbf{Respuesta correcta: C) 5 A}}

\textbf{Justificación:}\\
\textbf{Fórmula utilizada:} Resistencia equivalente en paralelo y corriente total (Lección: Sección 3 — Resistencias en paralelo):
$$R_{eq} = \frac{R_1 \cdot R_2}{R_1 + R_2} \qquad I_{total} = \frac{V}{R_{eq}}$$

\textbf{Sustitución:}
$$R_{eq} = \frac{20 \text{ }\Omega \times 30 \text{ }\Omega}{20 \text{ }\Omega + 30 \text{ }\Omega} = \frac{600}{50} = 12 \text{ }\Omega$$
$$I_{total} = \frac{60 \text{ V}}{12 \text{ }\Omega} = 5 \text{ A}$$

\textit{Verificación por ramas:} $I_1 = 60/20 = 3$ A; $I_2 = 60/30 = 2$ A; $I_{total} = 3 + 2 = 5$ A. \checkmark

\textbf{Respuestas incorrectas:}
\begin{itemize}
    \item \textbf{A) 2 A:} Corresponde solo a la corriente de la rama de 30 $\Omega$: $60/30 = 2$ A — se ignoró la otra rama.
    \item \textbf{B) 3 A:} Corresponde solo a la corriente de la rama de 20 $\Omega$: $60/20 = 3$ A — se ignoró la otra rama.
    \item \textbf{D) 12 A:} Error por dividir $V$ entre $R_{eq}$ calculado incorrectamente como $50/4 = 12.5$, o por confundir $R_{eq}$ con la suma directa y hacer $60/(20+30) = 1.2$ A y luego escalar mal.
\end{itemize}
\end{respuestabox}

\begin{respuestabox}
\subsection{Pregunta 7}
\textcolor{verdecorrecto}{\textbf{Respuesta correcta: D) 10 $\mu$F}}

\textbf{Justificación:}\\
\textbf{Fórmula utilizada:} Condensadores en paralelo (Lección: Sección 4 — Condensadores en serie y en paralelo):
$$C_{eq} = C_1 + C_2 + C_3$$

\textbf{Sustitución:}
$$C_{eq} = 3 \text{ }\mu\text{F} + 5 \text{ }\mu\text{F} + 2 \text{ }\mu\text{F} = 10 \text{ }\mu\text{F}$$

\textbf{Respuestas incorrectas:}
\begin{itemize}
    \item \textbf{A) 1 $\mu$F:} Error por aplicar la fórmula de serie ($1/C_{eq} = 1/3+1/5+1/2$) y calcular mal el resultado.
    \item \textbf{B) 3 $\mu$F:} Error por aplicar la fórmula de condensadores en serie: $1/C_{eq} = 1/3+1/5+1/2 \approx 1.03 \Rightarrow C_{eq} \approx 0.97\,\mu$F; distractor que confunde con otro resultado.
    \item \textbf{C) 8 $\mu$F:} Error por omitir uno de los tres condensadores: $3 + 5 = 8$ — olvida sumar el tercero de 2 $\mu$F.
\end{itemize}
\end{respuestabox}

\begin{respuestabox}
\subsection{Pregunta 8}
\textcolor{verdecorrecto}{\textbf{Respuesta correcta: C) 2 J}}

\textbf{Justificación:}\\
\textbf{Fórmula utilizada:} Energía almacenada en un condensador (Lección: Sección 4 — Capacitancia):
$$U = \frac{1}{2} C V^2$$

\textbf{Sustitución:}
$$U = \frac{1}{2} \times (100 \times 10^{-6} \text{ F}) \times (200 \text{ V})^2 = \frac{1}{2} \times 10^{-4} \text{ F} \times 40000 \text{ V}^2 = 2 \text{ J}$$

\textbf{Respuestas incorrectas:}
\begin{itemize}
    \item \textbf{A) 0.5 J:} Error por omitir el cuadrado del voltaje: $U = \frac{1}{2} C V = \frac{1}{2} \times 10^{-4} \times 200 = 0.01$ J — resultado inconsistente, distractor por confusión conceptual.
    \item \textbf{B) 1 J:} Error por olvidar el factor $\frac{1}{2}$: $C V^2 = 10^{-4} \times 4\times10^4 = 4$ J y luego dividir entre 2... o por usar $V = 100$ V en lugar de 200 V.
    \item \textbf{D) 20 J:} Error por no convertir $C$ a farads: usar $C = 100$ F en lugar de $100 \times 10^{-6}$ F; o por un error de potencia de diez.
\end{itemize}
\end{respuestabox}

\begin{respuestabox}
\subsection{Pregunta 9}
\textcolor{verdecorrecto}{\textbf{Respuesta correcta: B) $20 \times 10^{-6}$ T}}

\textbf{Justificación:}\\
\textbf{Fórmula utilizada:} Campo magnético de un conductor rectilíneo (Lección: Sección 5 — Campo magnético):
$$B = \frac{\mu_0 I}{2\pi r}$$

\textbf{Sustitución:}
$$B = \frac{(4\pi \times 10^{-7} \text{ T}\cdot\text{m/A})(10 \text{ A})}{2\pi (0.1 \text{ m})} = \frac{4\pi \times 10^{-6}}{2\pi} = \frac{4 \times 10^{-6}}{2} = 2 \times 10^{-5} \text{ T} = 20 \times 10^{-6} \text{ T}$$

\textbf{Respuestas incorrectas:}
\begin{itemize}
    \item \textbf{A) $10 \times 10^{-6}$ T:} Error por omitir el factor 2 en el denominador $2\pi r$, usando solo $\pi r$: $\mu_0 I / (\pi r) = 4\times10^{-6} / 0.1\pi$... o por dividir el resultado correcto entre 2 sin justificación.
    \item \textbf{C) $40 \times 10^{-6}$ T:} Error por duplicar el resultado, posiblemente por confundir $\mu_0 I / (\pi r)$ con $\mu_0 I / (2\pi r)$ y luego multiplicar por 2 de nuevo.
    \item \textbf{D) $200 \times 10^{-6}$ T:} Error por omitir completamente el $2\pi$ en el denominador: $\mu_0 I / r = 4\pi\times10^{-7}\times10/0.1 = 4\pi\times10^{-5} \approx 126\,\mu$T; o por un error de potencia de diez al operar con $r = 0.01$ m.
\end{itemize}
\end{respuestabox}

\begin{respuestabox}
\subsection{Pregunta 10}
\textcolor{verdecorrecto}{\textbf{Respuesta correcta: C) 100 V}}

\textbf{Justificación:}\\
\textbf{Fórmula utilizada:} Ley de Faraday (magnitud) (Lección: Sección 6 — Ley de Faraday):
$$|\mathcal{E}| = N \frac{|\Delta\Phi_B|}{\Delta t}$$

\textbf{Sustitución:}
$$|\mathcal{E}| = 200 \times \frac{0.05 \text{ Wb}}{0.1 \text{ s}} = 200 \times 0.5 \text{ V} = 100 \text{ V}$$

\textbf{Respuestas incorrectas:}
\begin{itemize}
    \item \textbf{A) 10 V:} Error por dividir $\Delta\Phi_B / \Delta t$ correctamente pero luego dividir entre $N$ en lugar de multiplicar: $0.5/200 \neq$ 10 — confusión en el despeje.
    \item \textbf{B) 50 V:} Error por usar $N = 100$ en lugar de $N = 200$: $100 \times 0.5 = 50$ V.
    \item \textbf{D) 1000 V:} Error por invertir el cociente de flujo: $\Delta t / \Delta\Phi_B = 0.1/0.05 = 2$ en lugar de 0.5, y luego $200 \times 2 \times 2.5 = $ valor inconsistente; o por usar $\Delta t = 0.01$ s.
\end{itemize}
\end{respuestabox}

\begin{respuestabox}
\subsection{Pregunta 11}
\textcolor{verdecorrecto}{\textbf{Respuesta correcta: C) 6 A}}

\textbf{Justificación:}\\
\textbf{Fórmulas utilizadas:} Paralelo y serie de resistencias, Ley de Ohm (Lección: Sección 3):

\textbf{Paso 1 — Resistencia equivalente del grupo en paralelo:}
$$R_{par} = \frac{12 \text{ }\Omega \times 12 \text{ }\Omega}{12 \text{ }\Omega + 12 \text{ }\Omega} = \frac{144}{24} = 6 \text{ }\Omega$$

\textbf{Paso 2 — Resistencia total (serie con $R_{par}$):}
$$R_{eq} = 6 \text{ }\Omega + 6 \text{ }\Omega = 12 \text{ }\Omega$$

\textbf{Paso 3 — Corriente total:}
$$I = \frac{V}{R_{eq}} = \frac{72 \text{ V}}{12 \text{ }\Omega} = 6 \text{ A}$$

\textbf{Respuestas incorrectas:}
\begin{itemize}
    \item \textbf{A) 3 A:} Error por sumar las tres resistencias en serie ($6+12+12 = 30\,\Omega$) ignorando que las dos de 12 $\Omega$ van en paralelo: $72/30 \neq 3$.
    \item \textbf{B) 4 A:} Error por calcular $R_{par} = 12/2 = 6\,\Omega$ correctamente pero luego usar $R_{eq} = 6+12 = 18\,\Omega$: $72/18 = 4$.
    \item \textbf{D) 12 A:} Error por ignorar la resistencia en serie y calcular solo con $R_{par}$: $72/6 = 12$.
\end{itemize}
\end{respuestabox}

\begin{respuestabox}
\subsection{Pregunta 12}
\textcolor{verdecorrecto}{\textbf{Respuesta correcta: C) 3 A}}

\textbf{Justificación:}\\
\textbf{Continuación del circuito de la Pregunta 11.}

\textbf{Paso 1 — Voltaje en el grupo paralelo:}\\
La corriente total es $I = 6$ A (calculada en P11). Esta corriente pasa por la resistencia de 6 $\Omega$ en serie y luego por el grupo paralelo ($R_{par} = 6\,\Omega$).
$$V_{par} = I \cdot R_{par} = 6 \text{ A} \times 6 \text{ }\Omega = 36 \text{ V}$$

\textbf{Paso 2 — Corriente en cada rama de 12 $\Omega$:}\\
En paralelo, cada rama tiene el mismo voltaje ($V_{par} = 36$ V):
$$I_{rama} = \frac{V_{par}}{R_{rama}} = \frac{36 \text{ V}}{12 \text{ }\Omega} = 3 \text{ A}$$

\textit{Verificación:} $I_{rama1} + I_{rama2} = 3 + 3 = 6$ A $= I_{total}$. \checkmark

\textbf{Respuestas incorrectas:}
\begin{itemize}
    \item \textbf{A) 1.5 A:} Error por dividir la corriente total entre las dos ramas: $I_{total}/2 = 6/2 = 3$ A es correcto en este caso simétrico, pero si se usa $I_{total} = 4$ A del error de P11, se obtiene $4/2 = 2$ A, no 1.5 A — distractor por cadena de error.
    \item \textbf{B) 2 A:} Error encadenado desde P11: si se usó $I_{total} = 4$ A, entonces $4/2 = 2$ A por rama.
    \item \textbf{D) 6 A:} Error por usar la corriente total directamente en cada rama, sin considerar que la corriente se divide en el paralelo.
\end{itemize}
\end{respuestabox}

\begin{respuestabox}
\subsection{Pregunta 13}
\textcolor{verdecorrecto}{\textbf{Respuesta correcta: B) 4 $\mu$F}}

\textbf{Justificación:}\\
\textbf{Fórmulas utilizadas:} Condensadores en paralelo y en serie (Lección: Sección 4):

\textbf{Paso 1 — Capacitancia equivalente del grupo en paralelo:}
$$C_{par} = C_2 + C_3 = 4 \text{ }\mu\text{F} + 8 \text{ }\mu\text{F} = 12 \text{ }\mu\text{F}$$

\textbf{Paso 2 — Capacitancia total ($C_1$ en serie con $C_{par}$):}
$$C_{eq} = \frac{C_1 \cdot C_{par}}{C_1 + C_{par}} = \frac{6 \text{ }\mu\text{F} \times 12 \text{ }\mu\text{F}}{6 \text{ }\mu\text{F} + 12 \text{ }\mu\text{F}} = \frac{72}{18} = 4 \text{ }\mu\text{F}$$

\textbf{Respuestas incorrectas:}
\begin{itemize}
    \item \textbf{A) 2 $\mu$F:} Error por aplicar serie a los tres condensadores directamente: $1/C_{eq} = 1/6+1/4+1/8 = 4/24+6/24+3/24 = 13/24 \Rightarrow C_{eq} \approx 1.85\,\mu$F, redondeado a 2.
    \item \textbf{C) 6 $\mu$F:} Error por confundir la configuración: sumar los dos en paralelo ($4+8=12$) y luego calcular en paralelo con 6 en lugar de en serie: $6+12=18$ y dividir entre algún factor.
    \item \textbf{D) 18 $\mu$F:} Error por sumar los tres condensadores directamente como si todos estuvieran en paralelo: $6+4+8 = 18$ — ignora la conexión en serie de $C_1$.
\end{itemize}
\end{respuestabox}

\begin{respuestabox}
\subsection{Pregunta 14}
\textcolor{verdecorrecto}{\textbf{Respuesta correcta: A) 0.1 Wb}}

\textbf{Justificación:}\\
\textbf{Fórmula utilizada:} Flujo magnético (Lección: Sección 6 — Flujo magnético):
$$\Phi_B = B \cdot A \cdot \cos\theta$$

\textbf{Atención al ángulo:} El problema indica que el \textit{plano} de la espira forma 30$^\circ$ con el campo. El ángulo $\theta$ en la fórmula se mide entre el campo y la \textit{normal} al plano, por lo que $\theta = 90^\circ - 30^\circ = 60^\circ$.

\textbf{Sustitución:}
$$\Phi_B = 0.5 \text{ T} \times 0.4 \text{ m}^2 \times \cos 60^\circ = 0.5 \times 0.4 \times 0.5 = 0.1 \text{ Wb}$$

\textbf{Respuestas incorrectas:}
\begin{itemize}
    \item \textbf{B) 0.17 Wb:} Error por usar $\theta = 30^\circ$ directamente: $0.5 \times 0.4 \times \cos 30^\circ = 0.2 \times (\sqrt{3}/2) \approx 0.17$ Wb — confusión entre el ángulo del plano y el ángulo de la normal.
    \item \textbf{C) 0.2 Wb:} Error por omitir el coseno: $\Phi_B = B \times A = 0.5 \times 0.4 = 0.2$ Wb — se usa $\cos\theta = 1$, como si el campo fuera perpendicular al plano.
    \item \textbf{D) 0.35 Wb:} Error por usar $\sin 30^\circ$ en lugar de $\cos 60^\circ$ y multiplicar por un factor adicional incorrecto.
\end{itemize}
\end{respuestabox}

\begin{respuestabox}
\subsection{Pregunta 15}
\textcolor{verdecorrecto}{\textbf{Respuesta correcta: C) 200 vueltas}}

\textbf{Justificación:}\\
\textbf{Fórmula utilizada:} Ley de Faraday despejada para $N$ (Lección: Sección 6 — Ley de Faraday):
$$|\mathcal{E}| = N \frac{|\Delta\Phi_B|}{\Delta t} \implies N = \frac{|\mathcal{E}| \cdot \Delta t}{|\Delta\Phi_B|}$$

\textbf{Sustitución:}
$$N = \frac{300 \text{ V} \times 0.04 \text{ s}}{0.06 \text{ Wb}} = \frac{12 \text{ V}\cdot\text{s}}{0.06 \text{ Wb}} = 200 \text{ vueltas}$$

\textbf{Respuestas incorrectas:}
\begin{itemize}
    \item \textbf{A) 50 vueltas:} Error por invertir $\Delta\Phi_B$ y $\Delta t$ en el despeje: $N = \mathcal{E} \cdot \Delta\Phi_B / \Delta t = 300 \times 0.06 / 0.04 = 450$, no 50 — distractor por error aritmético secundario.
    \item \textbf{B) 100 vueltas:} Error por usar $\Delta t = 0.02$ s (la mitad del valor real): $N = 300 \times 0.02 / 0.06 = 100$ — error de lectura del dato.
    \item \textbf{D) 450 vueltas:} Error por no despejar $N$ correctamente: se multiplica $\mathcal{E}$ por $\Delta\Phi_B$ y se divide entre $\Delta t$ en lugar de hacer lo inverso: $300 \times 0.06 / 0.04 \neq 450$ — confusión en el álgebra del despeje.
\end{itemize}
\end{respuestabox}
